\section{\kcode{dyn_groupprivate} Clause}
\label{sec:dyn_groupprivate}
\index{clauses!dyn_groupprivate@\kcode{dyn_groupprivate}}
\index{dyn_groupprivate clause@\kcode{dyn_groupprivate} clause}
\index{routines!omp_get_dyn_gprivate_ptr@\kcode{omp_get_dyn_gprivate_ptr}}
\index{omp_get_dyn_gprivate_ptr routine@\kcode{omp_get_dyn_gprivate_ptr} routine}


The following example shows the use of the \kcode{dyn_groupprivate} clause 
and the \kcode{omp_get_dyn_gprivate_ptr} routine. The \kcode{dyn_groupprivate} 
clause specifies the instantiation of dynamic \kcode{groupprivate} blocks in the \kcode{target teams} 
region. The block holds \ucode{N} integer values followed by \ucode{M} integer 
values, which total \ucode{N*sizeof(int)+M*sizeof(int)} bytes. Since the access-group 
modifier is not specified, it defaults to contention group access (\kcode{omp_access_cgroup}). 
Each contention group executing in the region will receive one of the uninitialized 
\kcode{groupprivate} buffer copies, which, in this particular example, means that each 
initiated team will effectively have its own copy.

The absolute size limit of a \kcode{groupprivate} memory space on a specific device 
can be obtained using the \kcode{omp_get_gprivate_limit} routine. The effective size 
that can be used in a \kcode{dyn_groupprivate} clause for a particular \kcode{target teams} region 
may be less than the absolute limit due to the use of the \kcode{groupprivate} directive 
or the OpenMP implementation. If the memory requested in the \kcode{dyn_groupprivate} 
clause exceeds the effective limit, the runtime will execute a fallback action. The fallback 
can be runtime error termination (abort), using an alternative implementation-defined 
memory space (\kcode{default_mem}), or continuing without instantiating the block (null). Since 
the fallback modifier is not specified, the fallback behavior defaults to \kcode{default_mem}.

The \kcode{omp_get_dyn_gprivate_ptr} routine returns a pointer to the dynamic 
\kcode{groupprivate} block for the current contention group (and \kcode{team}) in the 
\kcode{target teams} region, with an optional offset in bytes. The \kcode{omp_get_dyn_gprivate_ptr} calls used inside \ucode{foo} return pointers to the dynamic \kcode{groupprivate} block with 
offsets of \ucode{0} and \ucode{N*sizeof(int)} bytes, which correspond to the memory used to store 
the variables \ucode{x} and \ucode{y}, respectively. All threads belonging to the same contention group 
(and in this case, the same team) will have access to the same memory buffers, allowing 
the sharing of data within the same team. The lifetime of the copy is limited to the lifetime 
of all members of the corresponding contention group.


\cppexample[6.1]{dyn_groupprivate}{1}

\ffreeexample[6.1]{dyn_groupprivate}{1}
